% =====================================================================
%  LaTeX Workshop -- your working file
% =====================================================================
%
%  Everything on a line after a % is a comment: LaTeX ignores it.
%  That is what you are reading right now.
%
%  WHAT THIS DOCUMENT IS
%  A short report on how quickly a hot drink goes cold. It is a worked
%  example, not a real experiment -- nobody measured any coffee, and the
%  numbers and the plot were invented for this workshop. Your job today
%  is to TYPESET the report, not to run it.
%
%  By this evening it will contain a title, a table of contents,
%  sections, lists, equations, a table, a figure and a bibliography --
%  i.e. every ingredient of a term paper or a thesis.
%
%  HOUSE STYLE USED IN THESE FILES
%    * one sentence per line -- a long one may break at a comma
%    * \label goes on the same line as the \section it belongs to
%    * four spaces per indent level
%  None of this changes the output by a single pixel. It changes how
%  readable the SOURCE is, and how useful `git diff` is later.
%
%  IF YOU GET STUCK
%  Do not spend the rest of the workshop debugging. Open the matching
%  file in checkpoints/, copy its contents over this file, and carry on:
%
%    checkpoint-00.tex  this file, untouched
%    checkpoint-01.tex  state after block 1   (skeleton)
%    checkpoint-02.tex  state after block 2   (text and lists)
%    checkpoint-03.tex  state after block 3   (maths)
%    checkpoint-04.tex  the finished document
%
%  Each checkpoint already contains the TODOs for the block that follows
%  it, so you can always rejoin at the start of a block.
%
% =====================================================================
%  BLOCK 1 -- THE SKELETON
% =====================================================================
%
%  TODO 1.1  Write the smallest document that compiles. Below this
%            comment block, type:
%
%              \documentclass[a4paper,11pt]{article}
%
%              \begin{document}
%              Hello OST!
%              \end{document}
%
%            Then compile it. Look at the PDF. You have just typeset a
%            document.
%
%  TODO 1.2  Give it a title block. Above \begin{document}, add
%            \title{}, \author{} and \date{}; inside the document, as
%            the very first line, add \maketitle.
%            Title:  How Fast Does Your Coffee Go Cold?
%            Author: your own name
%            Date:   \today
%
%  TODO 1.3  Add the five sections of the report, in this order:
%
%              Introduction, Theory, Setup, Results, Conclusion
%
%            Write one sentence of your own under Introduction. The
%            other sections stay empty for now -- we fill them later.
%
%  TODO 1.4  Add \tableofcontents directly after \maketitle and
%            compile. Nobody typed that list of headings: it was built
%            from your \section commands, and it will stay correct when
%            you reorder or rename them.
%            This is the deal LaTeX offers you, and it is the single
%            biggest reason to put up with the syntax: you describe the
%            structure once, and the document maintains itself.
%
%  TODO 1.5  In your Introduction sentence, mark one term with \emph{},
%            one with \textbf{} and one short label with \texttt{}.
%
%  Stuck on 1.1? Checkpoint 01 has all of it, already done.
% =====================================================================

% Your document starts here:
